% =====================================================================
%  HomeGuard - Justifikasi keberadaan aplikasi & diferensiasi terhadap
%  alat sejenis yang sudah ada (menjawab pertanyaan penguji:
%  "sudah banyak aplikasi serupa").
%
%  Saran penempatan: sebagai sub-bab tambahan pada Bab II (Tinjauan
%  Pustaka), mis. "II.F Diferensiasi terhadap Alat yang Tersedia",
%  ATAU diintegrasikan ke akhir Bab II.E (Penelitian Terdahulu) dan
%  diperkuat di Bab I.A (Latar Belakang).
% =====================================================================

\subsection{Diferensiasi terhadap Alat yang Sudah Tersedia}

Patut diakui bahwa telah tersedia beragam perangkat pemindai jaringan,
baik komersial maupun \textit{open-source}. Namun, alat-alat tersebut
secara umum berada pada salah satu dari dua kutub yang keduanya tidak
sepenuhnya menjawab kebutuhan pengguna rumahan non-teknis dalam konteks
penilaian kerentanan IoT yang terstandarisasi.

Pada kutub \emph{perkakas profesional}, Nmap~\cite{b13}, Nessus, dan OpenVAS
merupakan pemindai yang sangat kuat namun ditujukan bagi praktisi
keamanan. Perkakas ini memiliki kurva belajar yang curam, keluarannya
bersifat teknis (daftar port/layanan mentah), berorientasi umum
(bukan spesifik IoT rumah), serta sebagian berlisensi atau memerlukan
infrastruktur dan konfigurasi yang tidak praktis bagi pengguna awam.
Yang terpenting, perkakas ini \emph{tidak} memetakan temuan ke taksonomi
kerentanan khusus IoT seperti OWASP IoT Top 10.

Pada kutub \emph{aplikasi konsumen}, perkakas seperti Fing dan
Bitdefender Home Scanner mudah digunakan, tetapi umumnya bersifat
\textit{closed-source} (basis aturannya tidak dapat diaudit), berbasis
\textit{cloud} sehingga metadata jaringan rumah dikirim ke server pihak
ketiga (menimbulkan persoalan privasi dan kedaulatan data), serta
menyajikan hasil sebagai label keamanan umum tanpa pemetaan eksplisit ke
kerangka akademik yang dapat dipertanggungjawabkan.

Dengan demikian, kontribusi HomeGuard \emph{bukan} terletak pada penemuan
teknik pemindaian baru --- HomeGuard memang memakai teknik standar
(\textit{TCP connect scan}, \textit{banner grabbing}, ARP/OUI
\textit{lookup}) --- melainkan pada \textbf{integrasi dan pemosisian}
yang mengisi celah spesifik berikut:

\begin{enumerate}
    \item \textbf{Pemetaan eksplisit ke OWASP IoT Top 10 (2018).} Setiap
    temuan diklasifikasikan ke kategori risiko IoT yang terstandarisasi
    dan disertai skor risiko serta rekomendasi mitigasi yang
    \textit{actionable}, bukan sekadar daftar port mentah.

    \item \textbf{Transparan dan dapat diaudit.} Sebagai \textit{open-source}
    dengan basis aturan terbuka, HomeGuard memungkinkan verifikasi
    ilmiah dan reproduktibilitas --- prasyarat bagi sebuah instrumen
    penelitian, yang tidak dimiliki perkakas \textit{closed-source}.

    \item \textbf{Privasi dan operasi lokal sepenuhnya.} Pemindaian dan
    analisis berjalan secara lokal tanpa mengirim data ke \textit{cloud},
    selaras dengan prinsip \textit{Zero Trust}~\cite{b19} dan menjaga kedaulatan
    data rumah tangga.

    \item \textbf{Ringan dan portabel.} Inti pemindaian hanya memakai
    pustaka standar Python tanpa dependensi wajib dan tanpa hak akses
    \textit{root}, sehingga dapat dijalankan oleh pengguna rumahan pada
    beragam lingkungan dengan mudah.

    \item \textbf{Berorientasi edukasi dan konteks lokal.} Antarmuka,
    catatan, dan rekomendasi disajikan dalam Bahasa Indonesia untuk
    meningkatkan kesadaran keamanan pengguna non-teknis, sekaligus
    memungkinkan studi empiris lanskap kerentanan IoT rumah tangga di
    Indonesia.
\end{enumerate}

Posisi pembeda ini diringkas pada Tabel~\ref{tab:diferensiasi}, yang
membandingkan HomeGuard dengan representasi alat pada kedua kutub di atas
berdasarkan kriteria yang relevan bagi pengguna rumahan.

\begin{table}[ht]
    \centering
    \caption{Diferensiasi HomeGuard terhadap Alat Sejenis}
    \label{tab:diferensiasi}
    \footnotesize
    \begin{tabular}{|p{2.2cm}|c|c|c|c|c|}
        \hline
        \textbf{Kriteria} & \textbf{Nmap} & \textbf{Nessus/} &
        \textbf{Fing} & \textbf{Bitdef.} & \textbf{Home-} \\
         & & \textbf{OpenVAS} & & \textbf{Scan} & \textbf{Guard} \\
        \hline
        \textit{Open-source} & Ya & Sebagian & Tidak & Tidak & \textbf{Ya} \\
        \hline
        Target pengguna awam & Tidak & Tidak & Ya & Ya & \textbf{Ya} \\
        \hline
        Pemetaan OWASP IoT & Tidak & Tidak & Tidak & Tidak & \textbf{Ya} \\
        \hline
        Lokal / tanpa cloud & Ya & Ya & Tidak & Tidak & \textbf{Ya} \\
        \hline
        Ringan / tanpa root & Tidak & Tidak & Ya & Ya & \textbf{Ya} \\
        \hline
        Edukasi (Bhs. lokal) & Tidak & Tidak & Tidak & Tidak & \textbf{Ya} \\
        \hline
    \end{tabular}
\end{table}

Singkatnya, tidak ada satu pun perkakas terdahulu yang secara simultan
bersifat \textit{open-source}, ramah pengguna awam, beroperasi lokal demi
privasi, ringan, sekaligus memetakan temuan ke OWASP IoT Top 10 dengan
rekomendasi edukatif. Kombinasi atribut inilah --- bukan teknik
pemindaiannya semata --- yang menjadi kebaruan dan justifikasi keberadaan
HomeGuard.
